\documentclass[border=7bp]{standalone}
% Shared preamble for the report's editable vector figures.
\usepackage{fontspec}
\usepackage{xeCJK}
\setmainfont{Calibri}
\setCJKmainfont{SimSun}[AutoFakeBold=1.5]
\usepackage{amsmath}
\usepackage{tikz}
\usetikzlibrary{arrows.meta,positioning,calc,shapes.geometric,fit}
\definecolor{Ink}{HTML}{25313D}
\definecolor{Blue}{HTML}{315E80}
\definecolor{PaleBlue}{HTML}{E7EFF5}
\definecolor{Muted}{HTML}{626B74}
\definecolor{PaleGray}{HTML}{F0F1F3}
\tikzset{
  >=Stealth,
  every node/.style={font=\fontsize{10bp}{13bp}\selectfont,text=Ink},
  flow/.style={->,draw=Ink,line width=.65pt},
  frame/.style={rectangle,draw=Ink,line width=.65pt,align=center,
    minimum height=.85cm,text width=5cm,inner sep=5pt},
  emphasis/.style={frame,draw=Blue,fill=PaleBlue},
  supporting/.style={frame,draw=Muted,fill=PaleGray},
  annotation/.style={font=\fontsize{9bp}{12bp}\selectfont,text=Muted,align=center},
}

\begin{document}
\begin{tikzpicture}
\path (-7.3,0) rectangle (7.3,0);

\node[emphasis,text width=5cm] (a) at (-3.2,0) {A*};
\node[emphasis,text width=5cm] (i) at (3.2,0) {BASIC IDA*};
\node[frame,text width=5cm] (a1) at (-3.2,-1.5) {Open：优先队列 + 状态索引\\保存待扩展的候选结点};
\node[frame,text width=5cm] (i1) at (3.2,-1.5) {当前递归路径 + 路径状态集合\\在 $f\leq B$ 内深度优先搜索};
\node[frame,text width=5cm] (a2) at (-3.2,-3.2) {Closed：已扩展状态\\Node 父链：用于恢复解路径};
\node[frame,text width=5cm] (i2) at (3.2,-3.2) {回溯后释放当前分支\\提高阈值后重新搜索};
\node[supporting,text width=5cm] (a3) at (-3.2,-4.9) {保存更多全局信息\\减少重复扩展，可能先耗尽内存};
\node[supporting,text width=5cm] (i3) at (3.2,-4.9) {路径相关存储随深度增长\\付出跨轮重复搜索的时间};
\node[frame,text width=11.5cm] (extra) at (0,-6.7) {额外存储分别计入：PDB 距离表；IDA* 可选的 TT / ZOBRIST；阈值诊断序列};
\foreach \a/\b in {a/a1,a1/a2,a2/a3,i/i1,i1/i2,i2/i3}{\draw[flow](\a)--(\b);}
\node[annotation,text width=11cm] at (0,-7.65) {框面积不表示内存大小；本实验未测 JVM 峰值内存，不把条目数或堆上限当作实际用量};

\end{tikzpicture}
\end{document}
